%\title{University of Bristol Thesis Template}
\RequirePackage[l2tabu]{nag}		% Warns for incorrect (obsolete) LaTeX usage
%
%
% File: memoirthesis.tex
% Author: Victor Brena
% Description: Contains the thesis template using memoir class,
% which is mainly based on book class but permits better control of
% chapter styles for example. This template is an adaptation and
% modification of Oscar's.
%
% Memoir is a flexible class for typesetting poetry, fiction,
% non-fiction and mathematical works as books, reports, articles or
% manuscripts. CTAN repository is found at:
% http://www.ctan.org/tex-archive/macros/latex/contrib/memoir/
%
%
% UoB guidelines for thesis presentation were found at:
% http://www.bris.ac.uk/esu/pg/pgrcop11-12topic.pdf#page=49
%
% UoB guidelines:
%
% The dissertation must be printed on A4 white paper. Paper up to A3 may be used
% for maps, plans, diagrams and illustrative material. Pages (apart from the
% preliminary pages) should normally be double-sided.
%
% Memoir class loads useful packages by default (see manual).
\documentclass[oneside,a4paper,12pt,reqno,openbib]{memoir} %add 'draft' to turn draft option on (see below)

% Adding metadata:
\usepackage{listings}
% SUMIT
\lstset{
basicstyle=\ttfamily,
columns=fullflexible,
breaklines=true,
postbreak=\mbox{\textcolor{red}{$\hookrightarrow$}\space},
}
% SUMIT

\usepackage{relsize}
\usepackage[table]{xcolor}
\usepackage{subcaption}
\usepackage{algorithm,algorithmicx,algpseudocode}
%\renewcommand{\algorithmcfname}{ALGORITHM}
%\SetAlFnt{\small}
%\SetAlCapFnt{\small}
%\SetAlCapNameFnt{\small}
%\SetAlCapHSkip{0pt}
%\IncMargin{-\parindent}

\setlength{\parskip}{1em}

\usepackage{datetime}
\usepackage{ifpdf}
\ifpdf
\pdfinfo{
   /Author (Sumit Kumar Yadav)
   /Title (Undergraduate Dissertation)
   /CreationDate (D:\pdfdate)
}
\fi
% When draft option is on.
\ifdraftdoc
	\usepackage{draftwatermark}				%Sets watermarks up.
	\SetWatermarkScale{0.3}
	\SetWatermarkText{\bf Draft: \today}
\fi
%
% Declare figure/table as a subfloat.
\newsubfloat{figure}
\newsubfloat{table}
% Better page layout for A4 paper, see memoir manual.
\settrimmedsize{297mm}{210mm}{*}
\setlength{\trimtop}{0pt}
\setlength{\trimedge}{\stockwidth}
\addtolength{\trimedge}{-\paperwidth}
\settypeblocksize{634pt}{448.13pt}{*}
\setulmargins{4cm}{*}{*}
\setlrmargins{*}{*}{0.6}
%\setlrmarginsandblock{1.5in}{1in}{*}
\setmarginnotes{17pt}{51pt}{\onelineskip}
\setheadfoot{\onelineskip}{2\onelineskip}
\setheaderspaces{*}{2\onelineskip}{*}
\checkandfixthelayout
%
\frenchspacing
% Font with math support: New Century Schoolbook
\usepackage{fouriernc}
\usepackage[T1]{fontenc}
%
% UoB guidelines:
%
% Text should be in double or 1.5 line spacing, and font size should be
% chosen to ensure clarity and legibility for the main text and for any
% quotations and footnotes. Margins should allow for eventual hard binding.
%
% Note: This is automatically set by memoir class. Nevertheless \OnehalfSpacing
% enables double spacing but leaves single spaced for captions for instance.
\OnehalfSpacing
%
% Sets numbering division level
\setsecnumdepth{subsection}
\maxsecnumdepth{subsubsection}
%
% Chapter style (taken and slightly modified from Lars Madsen Memoir Chapter
% Styles document
\usepackage{calc,soul,fourier}
\makeatletter
\newlength\dlf@normtxtw
\setlength\dlf@normtxtw{\textwidth}
\newsavebox{\feline@chapter}
\newcommand\feline@chapter@marker[1][4cm]{%
	\sbox\feline@chapter{%
		\resizebox{!}{#1}{\fboxsep=1pt%
			\colorbox{gray}{\color{white}\thechapter}%
		}}%
		\rotatebox{90}{%
			\resizebox{%
				\heightof{\usebox{\feline@chapter}}+\depthof{\usebox{\feline@chapter}}}%
			{!}{\scshape\so\@chapapp}}\quad%
		\raisebox{\depthof{\usebox{\feline@chapter}}}{\usebox{\feline@chapter}}%
}
\newcommand\feline@chm[1][4cm]{%
	\sbox\feline@chapter{\feline@chapter@marker[#1]}%
	\makebox[0pt][c]{% aka \rlap
		\makebox[1cm][r]{\usebox\feline@chapter}%
	}}
\makechapterstyle{daleifmodif}{
	\renewcommand\chapnamefont{\normalfont\Large\scshape\raggedleft\so}
	\renewcommand\chaptitlefont{\normalfont\Large\bfseries\scshape}
	\renewcommand\chapternamenum{} \renewcommand\printchaptername{}
	\renewcommand\printchapternum{\null\hfill\feline@chm[2.5cm]\par}
	\renewcommand\afterchapternum{\par\vskip\midchapskip}
	\renewcommand\printchaptertitle[1]{\color{black}\chaptitlefont\raggedleft ##1\par}
}
\renewcommand{\counterwithin}{\@ifstar{\@csinstar}{\@csin}}
\makeatother
\chapterstyle{daleifmodif}
%
% UoB guidelines:
%
% The pages should be numbered consecutively at the bottom centre of the
% page.
\makepagestyle{myvf}
\makeoddfoot{myvf}{}{\thepage}{}
\makeevenfoot{myvf}{}{\thepage}{}
\makeheadrule{myvf}{\textwidth}{\normalrulethickness}
\makeevenhead{myvf}{\small\textsc{\leftmark}}{}{}
\makeoddhead{myvf}{}{}{\small\textsc{\rightmark}}
\pagestyle{myvf}
%
% Oscar's command (it works):
% Fills blank pages until next odd-numbered page. Used to emulate single-sided
% frontmatter. This will work for title, abstract and declaration. Though the
% contents sections will each start on an odd-numbered page they will
% spill over onto the even-numbered pages if extending beyond one page
% (hopefully, this is ok).
\newcommand{\clearemptydoublepage}{\newpage{\thispagestyle{empty}\cleardoublepage}}
%
%
% Creates indexes for Table of Contents, List of Figures, List of Tables and Index
\makeindex
% \printglossaries below creates a list of abbreviations. \gls and related
% commands are then used throughout the text, so that latex can automatically
% keep track of which abbreviations have already been defined in the text.
%
% The import command enables each chapter tex file to use relative paths when
% accessing supplementary files. For example, to include
% chapters/brewing/images/figure1.png from chapters/brewing/brewing.tex we can
% use
% \includegraphics{images/figure1}
% instead of
% \includegraphics{chapters/brewing/images/figure1}
\usepackage{import}

% Add other packages needed for chapters here. For example:
\usepackage{lipsum}					%Needed to create dummy text
\usepackage{amsfonts} 					%Calls Amer. Math. Soc. (AMS) fonts
\usepackage[centertags]{amsmath}			%Writes maths centred down
\usepackage{stmaryrd}					%New AMS symbols
\usepackage{amssymb}					%Calls AMS symbols
\usepackage{amsthm}					%Calls AMS theorem environment
\usepackage{newlfont}					%Helpful package for fonts and symbols
\usepackage{mathtools}
\DeclarePairedDelimiter\abs{\lvert}{\rvert}%
\usepackage{layouts}					%Layout diagrams
\usepackage{graphicx}					%Calls figure environment
\graphicspath{{images/}}
\usepackage{wrapfig}
\usepackage{pdfpages}
\usepackage{lscape}
\usepackage{pdflscape}
\usepackage{rotating}
\usepackage{epstopdf}
%\usepackage{subfigure}
\usepackage{longtable,rotating}			%Long tab environments including rotation.
\usepackage[applemac]{inputenc}			%Needed to encode non-english characters
\usepackage{gensymb}					
\usepackage{colortbl}					%Makes coloured tables
\usepackage{wasysym}					%More math symbols
\usepackage{mathrsfs}					%Even more math symbols
\usepackage{float}						%Helps to place figures, tables, etc.
\usepackage{verbatim}					%Permits pre-formated text insertion
\usepackage{upgreek }					%Calls other kind of greek alphabet
\usepackage{latexsym}					%Extra symbols
\usepackage[square,numbers,
		     sort&compress]{natbib}		%Calls bibliography commands
\usepackage{url}						%Supports url commands
\usepackage{etex}						%eTeXÕs extended support for counters
\usepackage{fixltx2e}					%Eliminates some in felicities of the
									%original LaTeX kernel
\usepackage[spanish,english]{babel}		%For languages characters and hyphenation
\usepackage{color}                    				%Creates coloured text and background
\usepackage[colorlinks=true,
		     allcolors=black]{hyperref}              %Creates hyperlinks in cross references
\usepackage{memhfixc}					%Must be used on memoir document
									%class after hyperref
\usepackage{enumerate}					%For enumeration counter
\usepackage{footnote}					%For footnotes
\usepackage{microtype}					%Makes pdf look better.
\usepackage{rotfloat}					%For rotating and float environments as tables,
									%figures, etc.
\usepackage{alltt}						%LaTeX commands are not disabled in
									%verbatim-like environment
\usepackage[version=0.96]{pgf}			%PGF/TikZ is a tandem of languages for producing vector graphics from a
\usepackage{tikz}	
\usetikzlibrary{calc}
\usetikzlibrary{arrows,shapes,snakes,
		       automata,backgrounds,
		       petri,topaths}				%To use diverse features from tikz		
%							
%Reduce widows  (the last line of a paragraph at the start of a page) and orphans
% (the first line of paragraph at the end of a page)
\widowpenalty=1000
\clubpenalty=1000
%
% New command definitions for my thesis
%
\newcommand{\keywords}[1]{\par\noindent{\small{\bf Keywords:} #1}} %Defines keywords small section
\newcommand{\parcial}[2]{\frac{\partial#1}{\partial#2}}                             %Defines a partial operator
\newcommand{\vectorr}[1]{\mathbf{#1}}                                                        %Defines a bold vector
\newcommand{\vecol}[2]{\left(                                                                         %Defines a column vector
	\begin{array}{c}
		\displaystyle#1 \\
		\displaystyle#2
	\end{array}\right)}
\newcommand{\mados}[4]{\left(                                                                       %Defines a 2x2 matrix
	\begin{array}{cc}
		\displaystyle#1 &\displaystyle #2 \\
		\displaystyle#3 & \displaystyle#4
	\end{array}\right)}
\newcommand{\pgftextcircled}[1]{                                                                    %Defines encircled text
    \setbox0=\hbox{#1}%
    \dimen0\wd0%
    \divide\dimen0 by 2%
    \begin{tikzpicture}[baseline=(a.base)]%
        \useasboundingbox (-\the\dimen0,0pt) rectangle (\the\dimen0,1pt);
        \node[circle,draw,outer sep=0pt,inner sep=0.1ex] (a) {#1};
    \end{tikzpicture}
}
\newcommand{\range}[1]{\textnormal{range }#1}                                             %Defines range operator
\newcommand{\innerp}[2]{\left\langle#1,#2\right\rangle}                                 %Defines inner product
\newcommand{\prom}[1]{\left\langle#1\right\rangle}                                         %Defines average operator
\newcommand{\tra}[1]{\textnormal{tra} \: #1}                                                       %Defines trace operator
\newcommand{\sign}[1]{\textnormal{sign\,}#1}                                                   %Defines sign operator
\newcommand{\sech}[1]{\textnormal{sech} #1}                                                  %Defines sech
\newcommand{\diag}[1]{\textnormal{diag} #1}                                                    %Defines diag operator
\newcommand{\arcsech}[1]{\textnormal{arcsech} #1}                                       %Defines arcsech
\newcommand{\arctanh}[1]{\textnormal{arctanh} #1}                                         %Defines arctanh
%Change tombstone symbol
\newcommand{\blackged}{\hfill$\blacksquare$}
\newcommand{\whiteged}{\hfill$\square$}
\newcounter{proofcount}
\renewenvironment{proof}[1][\proofname.]{\par
 \ifnum \theproofcount>0 \pushQED{\whiteged} \else \pushQED{\blackged} \fi%
 \refstepcounter{proofcount}
 \normalfont
 \trivlist
 \item[\hskip\labelsep
       \itshape
   {\bf\em #1}]\ignorespaces
}{%
 \addtocounter{proofcount}{-1}
 \popQED\endtrivlist
}
%
%
% New definition of square root:
% it renames \sqrt as \oldsqrt
\let\oldsqrt\sqrt
% it defines the new \sqrt in terms of the old one
\def\sqrt{\mathpalette\DHLhksqrt}
\def\DHLhksqrt#1#2{%
\setbox0=\hbox{$#1\oldsqrt{#2\,}$}\dimen0=\ht0
\advance\dimen0-0.2\ht0
\setbox2=\hbox{\vrule height\ht0 depth -\dimen0}%
{\box0\lower0.4pt\box2}}
%
% My caption style
\newcommand{\mycaption}[2][\@empty]{
	\captionnamefont{\scshape}
	\changecaptionwidth
	\captionwidth{0.9\linewidth}
	\captiondelim{.\:}
	\indentcaption{0.75cm}
	\captionstyle[\centering]{}
	\setlength{\belowcaptionskip}{10pt}
	\ifx \@empty#1 \caption{#2}\else \caption[#1]{#2}
}

\newcommand*\vtick{\textsc{\char13}}
\newcommand*\expandedname{Twitter News Feed}
\newcommand*\shortname{TNF}

%
% My subcaption style
\newcommand{\mysubcaption}[2][\@empty]{
	\subcaptionsize{\small}
	\hangsubcaption
	\subcaptionlabelfont{\rmfamily}
	\sidecapstyle{\raggedright}
	\setlength{\belowcaptionskip}{10pt}
	\ifx \@empty#1 \subcaption{#2}\else \subcaption[#1]{#2}
}
%
%An initial of the very first character of the content
\usepackage{lettrine}
\newcommand{\initial}[1]{%
	\lettrine[lines=3,lhang=0.33,nindent=0em]{
		\color{gray}
     		{\textsc{#1}}}{}}
%
% Theorem styles used in my thesis
%
\theoremstyle{plain}
\newtheorem{theo}{Theorem}[chapter]
\theoremstyle{plain}
\newtheorem{prop}{Proposition}[chapter]
\theoremstyle{plain}
\theoremstyle{definition}
\newtheorem{dfn}{Definition}[chapter]
\theoremstyle{plain}
\newtheorem{lema}{Lemma}[chapter]
\theoremstyle{plain}
\newtheorem{cor}{Corollary}[chapter]
\theoremstyle{plain}
\newtheorem{resu}{Result}[chapter]
%
% Hyphenation for some words
%
\hyphenation{res-pec-tively}
\hyphenation{mono-ti-ca-lly}
\hyphenation{hypo-the-sis}
\hyphenation{para-me-ters}
\hyphenation{sol-va-bi-li-ty}
%
%
\begin{document}
% UoB guidlines:
%
% Preliminary pages
%
% The five preliminary pages must be the Title Page, Abstract, Dedication
% and Acknowledgements, Author's Declaration and Table of Contents.
% These should be single-sided.
%
% Table of contents, list of tables and illustrative material
%
% The table of contents must list, with page numbers, all chapters,
 % sections and subsections, the list of references, bibliography, list of
% abbreviations and appendices. The list of tables and illustrations
% should follow the table of contents, listing with page numbers the
% tables, photographs, diagrams, etc., in the order in which they appear
% in the text.
%
\frontmatter
\pagenumbering{roman}
%
\begin{titlingpage}
\begin{SingleSpace}
\calccentering{\unitlength}
\begin{adjustwidth*}{\unitlength}{-\unitlength}
\vspace*{-1.5 cm}
\begin{center}
{\LARGE \textbf{Architectural level design of an Island-Style FPGA}}\\
\vspace{1mm}
{\Large \textbf{}}\\[4mm]
\vspace{0.5 cm}
{\large PROJECT REPORT\\ \emph{submitted towards the partial fulfillment of the\\ \vspace{1 mm} requirements for the award of the degree \\ \vspace{1 mm} of}}\\
\vspace{5 mm}
{\large \textbf{BACHELOR OF TECHNOLOGY} \vspace{3 mm} \\\emph{in}\\\vspace{3 mm} \textbf{ELECTRONICS AND COMMUNICATION ENGINEERING}}\\
\vspace{4 mm}
%{\large \textbf{(With specialization in Information Technology)}} \\
\vspace{2 mm}
{\large \emph{Submitted by}}\\
\vspace{5mm}

\begin{minipage}[t]{0.4\textwidth}
\begin{center} \large
{\large\textbf{SUMIT KUMAR YADAV}}\\
{15117068} % Author name - remove the \href bracket to remove the link
\end{center}
\end{minipage}

%---------------------------------------
\vspace{0 mm}
{\large \textit{Under the guidance of:}}\\
\vspace{3 mm}
{\large\textbf{Dr. ANAND BULUSU}}\\%---------------------------------------
\vspace{8 mm}
\includegraphics[height = 4cm, keepaspectratio]{IITR_Logo.jpg}\\
\vspace{8 mm}
{\large \textbf{DEPARTMENT OF ELECTRONICS AND COMMUNICATION ENGINEERING}\\}
\vspace{1mm}
{\large \textbf{INDIAN INSTITUTE OF TECHNOLOGY ROORKEE}}\\
\vspace{1mm}
{\large \textbf{ROORKEE $-$ 247667 (INDIA)}}\\
\vspace{1mm}
{\large \textbf{May, 2019}}
%\vspace{5 mm}
\end{center}
%\begin{flushright}
%{\small Word count: ten thousand and four}
%\end{flushright}
\end{adjustwidth*}
\end{SingleSpace}
\end{titlingpage} 

\clearemptydoublepage
%
%
% File: declaration.tex
% Author: V?ctor Bre?a-Medina
% Description: Contains the declaration page
%
% UoB guidelines:
%
% Author's declaration
%
% I declare that the work in this dissertation was carried out in accordance
% with the requirements of the University's Regulations and Code of Practice
% for Research Degree Programmes and that it has not been submitted for any
% other academic award. Except where indicated by specific reference in the
% text, the work is the candidate's own work. Work done in collaboration with,
% or with the assistance of, others, is indicated as such. Any views expressed
% in the dissertation are those of the author.
%
% SIGNED: .............................................................
% DATE:..........................
%
%\chapter*{Candidate's Declaration}
\noindent \hfill \textsc{\Large Student's Declaration}\\
\noindent\makebox[\linewidth]{\rule{\textwidth}{1pt}} 
\begin{SingleSpacing}
%\begin{quote}
\vspace{-5mm}
I declare that the work presented in this report with title ``\textbf{Architectural level design of an Island-Style FPGA}" towards the fulfillment of the requirement for the award of the degree of \textbf{Bachelor of Technology} in \textbf{Electronics \& Communication Engineering} submitted in the \textbf{Department of Electronics \& Communication Engineering}, \textbf{Indian Institute of Technology, Roorkee}, India is an authentic record of my own work carried out during the period \textbf{from August 2018 to May 2019} under the supervision of \textbf{Dr. Anand Bulusu}, Associate Professor, Dept. of ECE, IIT Roorkee.
\noindent The content of this report has not been submitted by me for the award of any other degree of this or any other institute.

\vspace{3cm}
\noindent
\textsc{DATE : ......................... ~~~~~~~~~~~~~~~~~~~~~~~~~~~~~~~~~~SIGNED: ........................................}\\\\
\textsc{PLACE: .........................}
\nobreak
\hfill \textsc{(SUMIT KUMAR YADAV)} \\
\vspace{-4 mm}
\nobreak
\hfill \textsc{Enrol. No.: 15117068}


\end{SingleSpacing}
\clearpage 

\clearemptydoublepage


%
% File: declaration.tex
% Author: V?ctor Bre?a-Medina
% Description: Contains the declaration page
%
% UoB guidelines:
%
% Author's declaration
%
% I declare that the work in this dissertation was carried out in accordance
% with the requirements of the University's Regulations and Code of Practice
% for Research Degree Programmes and that it has not been submitted for any
% other academic award. Except where indicated by specific reference in the
% text, the work is the candidate's own work. Work done in collaboration with,
% or with the assistance of, others, is indicated as such. Any views expressed
% in the dissertation are those of the author.
%
% SIGNED: .............................................................
% DATE:..........................
%
\chapter*{Supervisor's Certificate}
\begin{OnehalfSpacing}
%\begin{quote}
This is to certify that the statement made by the candidate is correct to the best of my knowledge and belief.

\vspace{2.25cm}
\noindent
\textsc{DATE : .........................~~~~~~~~~~~~~~~~~~~~~~~~~~~~~~~~~~~SIGNED: ........................................}

\nobreak
\hfill \textsc{(Dr. ANAND BULUSU)}

\vspace{-4 mm}
\nobreak
\hfill \textsc{Associate Professor}

\vspace{-4 mm}
\nobreak
\hfill \textsc{Dept. of ECE, IIT Roorkee}
\end{OnehalfSpacing}
\clearpage 
\clearemptydoublepage
%

%
% file: dedication.tex
% author: V?ctor Bre?a-Medina
% description: Contains the text for thesis dedication
%

\chapter*{Acknowledgements}
\begin{OnehalfSpacing}
%\begin{center}
%\begin{minipage}{14.7 cm}
\initial I sincerely thank my B.Tech. Project advisor \textbf{Dr. Anand Bulusu} for his constant guidance and support throughout the project. 

I would like to acknowledge the contribution of \textbf{Mr. Jaskirat Singh} \& \textbf{Mr. Sanjay Yadav}, B. Tech students in ECE, IIT Roorkee for helping out with the necessary layouts and the ideation of the architecture. I am sincerely thankful to \textbf{Mr. Inder Choudhary}, research scholar in Microelectronics \& VLSI, Department of ECE, IIT Roorkee for his ideas and motivation. I also thank all the research scholars and project staff of SMDP Lab, Department of ECE for helpful technical discussions at every stage of the project. I thank \textbf{Department of Electronics and Communication Engineering, IIT Roorkee} for providing the lab and other resources required for carrying out my project work. 

Overall, working on this project has been a knowledgeable experience and provided me with a great experience in the area of FPGAs \& computer architecture. 


%\end{minipage}
%\end{center}
\end{OnehalfSpacing}
\clearpage

\clearemptydoublepage
%
%
% File: abstract.tex
% Author: V?ctor Bre?a-Medina
% Description: Contains the text for thesis abstract
%
% UoB guidelines:
%
% Each copy must include an abstract or summary of the dissertation in not
% more than 300 words, on one side of A4, which should be single-spaced in a
% font size in the range 10 to 12. If the dissertation is in a language other
% than English, an abstract in that language and an abstract in English must
% be included.

\chapter*{Abstract}
%\begin{SingleSpace}
\begin{OnehalfSpacing}
\paragraph{}
\end{OnehalfSpacing}
Minimal FPGAs have huge promise for implementing small interfacing and other logic at cheaper costs. The minimal FPGAs available in the market are still too big for certain applications and thus we try to solve this problem. We developed the flow for designing a minimal Island-Style FPGA using VTR[3] experimentation and Cadence Schematic level simulations. We first develop models for area and delay of a CLB and choose suitable LUT size. After that, we experiment with different switch matrix configurations and routing widths to settle for the optimum. The I/O block for the FPGA is then designed in Cadence Schematic L and tested for functionality. SKILL language is used to automatically generate decoder schematics. All smaller modules are used to create the top level schematic of the FPGA. A flow is created to program the bitstream using stimulus files and some scripts. The FPGA is tested for functionality and the delay path is then analyzed. The breakdown of the delay is provided for reference. All the simulations and tests were performed in Cadence ADE L using spectre at room temperature. SCL $180nm$ PDK was used in tt18 configuration with minimal sizing unless explicitly stated. The thesis ends by concluding the results and discussing future directions.

\clearpage 

\clearemptydoublepage
%
\renewcommand{\contentsname}{Table of Contents}
\maxtocdepth{subsection}
\tableofcontents*
\addtocontents{toc}{\par\nobreak \mbox{}\hfill{\bf Page}\par\nobreak}
\clearemptydoublepage
%
%\listoftables
%\addtocontents{lot}{\par\nobreak\textbf{{\scshape Table} \hfill Page}\par\nobreak}
%\clearemptydoublepage
%
\listoffigures
\addtocontents{lof}{\par\nobreak\textbf{{\scshape Figure} \hfill Page}\par\nobreak}
\clearemptydoublepage
%
% The bulk of the document is delegated to these chapter files in
% subdirectories.
\mainmatter
%
\import{chapters/chapter01/}{chap01.tex}
\clearemptydoublepage
%
\import{chapters/chapter02/}{chap02.tex}
\clearemptydoublepage
%
\import{chapters/chapter03/}{chap03.tex}
\clearemptydoublepage
%
\import{chapters/chapter04/}{chap04.tex}
\clearemptydoublepage
%
\import{chapters/chapter05/}{chap05.tex}
\clearemptydoublepage
%
\import{chapters/chapter06/}{chap06.tex}
\clearemptydoublepage
%
\import{chapters/chapter07/}{chap07.tex}
\clearemptydoublepage
%
\import{chapters/chapter08/}{chap08.tex}
\clearemptydoublepage
%

%
%
%
% Apparently the guidelines don't say anything about citations or
% bibliography styles so I guess we can use anything.
\backmatter
\bibliographystyle{ieeetr}
\refstepcounter{chapter}
\bibliography{thesisbiblio}
\clearemptydoublepage

\appendix
\def\code#1{\texttt{#1}}
\chapter{Appendix A}
\label{chap:intro}
\vspace{-0.3 cm}


\section{A.I \qquad Architecture XML file of our FPGA}

\paragraph{}
\begin{lstlisting}
<architecture>

  <!-- 
       ODIN II specific config begins 
       Describes the types of user-specified netlist blocks (in blif, this corresponds to 
       ".model [type_of_block]") that this architecture supports.

       Note: Basic LUTs, I/Os, and flip-flops are not included here as there are 
       already special structures in blif (.names, .input, .output, and .latch) 
       that describe them.
  -->
  <models>
  </models>
  <!-- ODIN II specific config ends -->

  <!-- Physical descriptions begin -->
  <layout>
     <auto_layout aspect_ratio="1.000000">
        <!--Perimeter of 'io' blocks with 'EMPTY' blocks at corners-->
        <perimeter type="io" priority="100"/>
        <corners type="EMPTY" priority="101"/>
        <!--Fill with 'clb'-->
        <fill type="clb" priority="10"/>
        </auto_layout>
</layout>
	<device>
		<sizing R_minW_nmos="6065.520020" R_minW_pmos="18138.500000"/>
		<area grid_logic_tile_area="7238.080078"/>
		<chan_width_distr>
		<x distr="uniform" peak="1.000000"/>
			<y distr="uniform" peak="1.000000"/>
		</chan_width_distr>
		<switch_block type="wilton" fs="3"/>
	<connection_block input_switch_name="ipin_cblock"/>
        </device>
	<switchlist>
		<switch type="mux" name="0" R="0.000000" Cin="0.000000e+00" Cout="0.000000e+00" Tdel="7.958000e-11" mux_trans_size="2.074780" buf_size="19.261999"/>
	<!--switch ipin_cblock resistance set to yeild for 4x minimum drive strength buffer-->
        <switch type="mux" name="ipin_cblock" R="1516.380005" Cout="0." Cin="0.000000e+00" Tdel="7.362000e-11" mux_trans_size="1.240240" buf_size="auto"/>
        </switchlist>
	<segmentlist>
		<segment freq="1.000000" length="1" type="unidir" Rmetal="0.000000" Cmetal="0.000000e+00">
			<mux name="0"/>
			<sb type="pattern">1 1</sb>
			<cb type="pattern">1</cb>
		</segment>
	</segmentlist>

  <complexblocklist>

    <!-- Define I/O pads begin -->
    <!-- Capacity is a unique property of I/Os, it is the maximum number of I/Os that can be placed at the same (X,Y) location on the FPGA -->
    <pb_type name="io" capacity="1">
      <input name="outpad" num_pins="1"/>
      <output name="inpad" num_pins="1"/>
      <clock name="clock" num_pins="1"/>

      <!-- IOs can operate as either inputs or outputs.
	     Delays below come from Ian Kuon. They are small, so they should be interpreted as
	     the delays to and from registers in the I/O (and generally I/Os are registered 
	     today and that is when you timing analyze them.
	     -->
      <mode name="inpad">
        <pb_type name="inpad" blif_model=".input" num_pb="1">
          <output name="inpad" num_pins="1"/>
        </pb_type>
        <interconnect>
          <direct name="inpad" input="inpad.inpad" output="io.inpad">
            <delay_constant max="4.791000e-11" in_port="inpad.inpad" out_port="io.inpad"/>
          </direct>
        </interconnect>

      </mode>
      <mode name="outpad">
        <pb_type name="outpad" blif_model=".output" num_pb="1">
          <input name="outpad" num_pins="1"/>
        </pb_type>
        <interconnect>
          <direct name="outpad" input="io.outpad" output="outpad.outpad">
            <delay_constant max="1.557000e-11" in_port="io.outpad" out_port="outpad.outpad"/>
          </direct>
        </interconnect>
      </mode>

      <!-- Every input pin is driven by 100% of the tracks in a channel, every output pin is driven by 100% of the tracks in a channel -->
      <fc in_type="frac" in_val="1" out_type="frac" out_val="1"/>

      <!-- IOs go on the periphery of the FPGA, for consistency, 
          make it physically equivalent on all sides so that only one definition of I/Os is needed.
          If I do not make a physically equivalent definition, then I need to define 4 different I/Os, one for each side of the FPGA
        -->
      <pinlocations pattern="custom">
        <loc side="left">io.outpad io.inpad io.clock</loc>
        <loc side="top">io.outpad io.inpad io.clock</loc>
        <loc side="right">io.outpad io.inpad io.clock</loc>
        <loc side="bottom">io.outpad io.inpad io.clock</loc>
      </pinlocations>

      <!-- Place I/Os on the sides of the FPGA -->
      <power method="ignore"/>			
    </pb_type>
    <!-- Define I/O pads ends -->

    <!-- Define general purpose logic block (CLB) begin -->
    <pb_type name="clb">
      <input name="I" num_pins="4" equivalent="full"/>
      <output name="O" num_pins="1" equivalent="instance"/>
      <clock name="clk" num_pins="1"/>

      <!-- Describe basic logic element. -->
      <!-- Define 6-LUT mode -->
          <pb_type name="ble4" num_pb="1">
            <input name="in" num_pins="4"/>
            <output name="out" num_pins="1"/>
            <clock name="clk" num_pins="1"/> 

            <!-- Define LUT -->
            <pb_type name="lut4" blif_model=".names" num_pb="1" class="lut">
              <input name="in" num_pins="4" port_class="lut_in"/>
              <output name="out" num_pins="1" port_class="lut_out"/>
              <!-- LUT timing using delay matrix -->
              <delay_matrix type="max" in_port="lut4.in" out_port="lut4.out">
                2.063000e-10
                2.063000e-10
                2.063000e-10
                2.063000e-10
              </delay_matrix>
            </pb_type>

            <!-- Define flip-flop -->
            <pb_type name="ff" blif_model=".latch" num_pb="1" class="flipflop">
              <input name="D" num_pins="1" port_class="D"/>
              <output name="Q" num_pins="1" port_class="Q"/>
              <clock name="clk" num_pins="1" port_class="clock"/>
			  <!-- setup time included in LUT delay -->
              <T_setup value="0.000000e-10" port="ff.D" clock="clk"/>
              <T_clock_to_Q max="8.406000e-11" port="ff.Q" clock="clk"/>
            </pb_type>

            <interconnect>
              <direct name="direct1" input="ble4.in" output="lut4[0:0].in"/>
              <direct name="direct2" input="lut4.out" output="ff.D">
                <!-- Advanced user option that tells CAD tool to find LUT+FF pairs in netlist -->
                <pack_pattern name="ble6" in_port="lut4.out" out_port="ff.D"/>
              </direct>
              <direct name="direct3" input="ble4.clk" output="ff.clk"/>                    
              <mux name="mux1" input="ff.Q lut4.out" output="ble4.out">
              </mux>
            </interconnect>
          </pb_type>
      <interconnect>
        <!-- We use a full crossbar to get logical equivalence at inputs of CLB  -->
        <complete name="crossbar" input="clb.I ble4[0:0].out" output="ble4[0:0].in">
          <delay_constant max="5.043000e-11" in_port="clb.I" out_port="ble4[0:0].in"/>
          <delay_constant max="5.031000e-11" in_port="ble4[0:0].out" out_port="ble4[0:0].in"/>
        </complete>
        <complete name="clks" input="clb.clk" output="ble4[0:0].clk">
        </complete>
        <direct name="clbouts1" input="ble4[0:0].out" output="clb.O"/>
      </interconnect>

      <!-- Every input pin is driven by 100% of the tracks in a channel, every output pin is driven by 100% of the tracks in a channel -->
      <fc in_type="frac" in_val="1" out_type="frac" out_val="1"/>

      <pinlocations pattern="custom">
          <loc side="left">clb.I[0] clb.clk</loc> 
          <loc side="top">clb.I[1]</loc> 
          <loc side="right">clb.I[2] clb.O</loc> 
          <loc side="bottom">clb.I[3]</loc>     
        </pinlocations>

      <!-- Place this general purpose logic block in any unspecified column -->
      </pb_type>
    <!-- Define general purpose logic block (CLB) ends -->

  </complexblocklist>
</architecture>
\end{lstlisting}


\section{A.II \qquad Bitstream to binary stimulus - script}

\paragraph{}
\begin{lstlisting}
import sys
import csv
sys.stdout = open('memory_data_compact.txt', 'w')
data = list(csv.reader(open('memory_mapping.csv')))

sys.stdout.write('BL1       BL1       0 000 ')

for j in range(1,17):
	for k in range(1,10):
		sys.stdout.write(str(data[j][4*k-3]))
print(' 00000000000000')

sys.stdout.write('BL2       BL2       0 000 ')
for j in range(1,17):
	for k in range(1,10):
		sys.stdout.write(str(data[j][4*k-2]))
print(' 00000000000000')

sys.stdout.write('BL3       BL3       0 000 ')
for j in range(1,17):
	for k in range(1,10):
		sys.stdout.write(str(data[j][4*k-1]))
print(' 00000000000000')

sys.stdout.write('BL4       BL4       0 000 ')
for j in range(1,17):
	for k in range(1,10):
		sys.stdout.write(str(data[j][4*k]))
print(' 00000000000000')

\end{lstlisting}

\section{A.III \qquad binary sequences to .scs - script }

\paragraph{}
\begin{lstlisting}
import sys
import fileinput
orig_stdout = sys.stdout
sys.stdout = open('stimulus.scs', 'w')
TIM=147
print('simulator lang=spectre') 
print('parameters VDD = 1.8')
print('//tr is (transition time)/2')
print('parameters tr = 1n') 
print('//f is the fall time') 
print('//parameters f = 100p')
print('//K is the signal period for programming') 
print('parameters K = 5u') 
print('//L is the signal period for running') 
print('parameters L = 0.5u')
print('\n')
print('vVdd (vdd 0) vsource dc=VDD')
print('vGnd (gnd 0) vsource dc=0')
print('\n')

lines = [line.rstrip('\n') for line in open('memory_data_compact.txt','r')]

while '' in lines:
    lines.remove('')

for line in lines:
	words=line.split()
	print('v'+words[0]+' ('+words[1]+ ' '+words[2]+') vsource type=pwl wave=\[')
	sent=words[3:]
	sen=''.join(sent)
	sen=sen.replace(" ", "")
	bits=list(sen)
	for i in range(len(bits)):
		ostr=""
		ostr+="+ ("
		if i==0:
			ostr+=" 0*K   ) "
		else:
			if i<TIM:
				ostr+=str(i).rjust(2)+'*K+tr) '
			else:
				ostr+=str(TIM).rjust(2)+'*K+'+str(i-TIM).rjust(2)+'*L+tr) '
		if bits[i]=='0':
			if i<TIM:
				ostr+='0    ('+str(i+1).rjust(2)+'*K-tr) 0'
			else:
				ostr+='0    ('+str(TIM).rjust(2)+'*K+'+str(i-TIM+1).rjust(2)+'*L-tr) 0'
		if bits[i]=='1':
			if i<TIM:
				ostr+='VDD  ('+str(i+1).rjust(2)+'*K-tr) VDD'
			else:
				ostr+='VDD  ('+str(TIM).rjust(2)+'*K+'+str(i-TIM+1).rjust(2)+'*L-tr) VDD'
		print(ostr)
	print('+ ]')

	if words[0][0]=='B' and words[0][1]=='L':
		print('vBL_'+words[0][2:]+' (BL_'+words[1][2:]+ ' '+words[2]+') vsource type=pwl wave=\[')
		sent=words[3:]
		sen=''.join(sent)
		sen=sen.replace(" ", "")
		bits=list(sen)
		for i in range(len(bits)):
			ostr=""
			ostr+="+ ("
			if i==0:
				ostr+=" 0*K   ) "
			else:
				if i<TIM:
					ostr+=str(i).rjust(2)+'*K+tr) '
				else:
					ostr+=str(TIM).rjust(2)+'*K+'+str(i-TIM).rjust(2)+'*L+tr) '
			if bits[i]=='0':
				if i<TIM:
					ostr+='VDD  ('+str(i+1).rjust(2)+'*K-tr) VDD'
				else:
					ostr+='VDD  ('+str(TIM).rjust(2)+'*K+'+str(i-TIM+1).rjust(2)+'*L-tr) VDD'
			if bits[i]=='1':
				if i<TIM:
					ostr+='0    ('+str(i+1).rjust(2)+'*K-tr) 0'
				else:
					ostr+='0    ('+str(TIM).rjust(2)+'*K+'+str(i-TIM+1).rjust(2)+'*L-tr) 0'
			print(ostr)
		print('+ ]')
		
sys.stdout.close()
sys.stdout=orig_stdout 
# Read in the file
with open('stimulus.scs', 'r') as file :
  filedata = file.read()

# Replace the target string
filedata = filedata.replace('BL', 'BitLine')

# Write the file out again
with open('stimulus.scs', 'w') as file:
  file.write(filedata)

\end{lstlisting}

\section{A.IV \qquad .scs for D Flip Flop}

\paragraph{}
\begin{lstlisting}
simulator lang=spectre
parameters VDD = 1.8
//tr is (transition time)/2
parameters tr = 50p
//f is the fall time
//parameters f = 100p
//K is the signal period
parameters K = 40n


vVdd (vdd 0) vsource dc=VDD
vGnd (gnd 0) vsource dc=0


vPRE (PRE 0) vsource type=pwl wave=\[
+ ( 0*K   ) VDD  ( 1*K-tr) VDD
+ ( 1*K+tr) VDD  ( 2*K-tr) VDD
+ ( 2*K+tr) VDD  ( 3*K-tr) VDD
+ ( 3*K+tr) VDD  ( 4*K-tr) VDD
+ ( 4*K+tr) VDD  ( 5*K-tr) VDD
+ ( 5*K+tr) VDD  ( 6*K-tr) VDD
+ ( 6*K+tr) VDD  ( 7*K-tr) VDD
+ ( 7*K+tr) VDD  ( 8*K-tr) VDD
+ ( 8*K+tr) VDD  ( 9*K-tr) VDD
+ ( 9*K+tr) VDD  (10*K-tr) VDD
+ (10*K+tr) VDD  (11*K-tr) VDD
+ (11*K+tr) 0    (12*K-tr) 0
+ (12*K+tr) 0    (13*K-tr) 0
+ (13*K+tr) 0    (14*K-tr) 0
+ (14*K+tr) 0    (15*K-tr) 0
+ (15*K+tr) VDD  (16*K-tr) VDD
+ (16*K+tr) VDD  (17*K-tr) VDD
+ (17*K+tr) VDD  (18*K-tr) VDD
+ (18*K+tr) VDD  (19*K-tr) VDD
+ (19*K+tr) VDD  (20*K-tr) VDD
+ (20*K+tr) VDD  (21*K-tr) VDD
+ (21*K+tr) VDD  (22*K-tr) VDD
+ (22*K+tr) VDD  (23*K-tr) VDD
+ (23*K+tr) VDD  (24*K-tr) VDD
+ (24*K+tr) VDD  (25*K-tr) VDD
+ (25*K+tr) VDD  (26*K-tr) VDD
+ (26*K+tr) VDD  (27*K-tr) VDD
+ (27*K+tr) VDD  (28*K-tr) VDD
+ (28*K+tr) VDD  (29*K-tr) VDD
+ (29*K+tr) VDD  (30*K-tr) VDD
+ (30*K+tr) VDD  (31*K-tr) VDD
+ (31*K+tr) VDD  (32*K-tr) VDD
+ (32*K+tr) VDD  (33*K-tr) VDD
+ ]
vCLR (CLR 0) vsource type=pwl wave=\[
+ ( 0*K   ) 0    ( 1*K-tr) 0
+ ( 1*K+tr) 0    ( 2*K-tr) 0
+ ( 2*K+tr) 0    ( 3*K-tr) 0
+ ( 3*K+tr) 0    ( 4*K-tr) 0
+ ( 4*K+tr) 0    ( 5*K-tr) 0
+ ( 5*K+tr) 0    ( 6*K-tr) 0
+ ( 6*K+tr) 0    ( 7*K-tr) 0
+ ( 7*K+tr) 0    ( 8*K-tr) 0
+ ( 8*K+tr) 0    ( 9*K-tr) 0
+ ( 9*K+tr) 0    (10*K-tr) 0
+ (10*K+tr) 0    (11*K-tr) 0
+ (11*K+tr) VDD  (12*K-tr) VDD
+ (12*K+tr) VDD  (13*K-tr) VDD
+ (13*K+tr) VDD  (14*K-tr) VDD
+ (14*K+tr) VDD  (15*K-tr) VDD
+ (15*K+tr) VDD  (16*K-tr) VDD
+ (16*K+tr) VDD  (17*K-tr) VDD
+ (17*K+tr) VDD  (18*K-tr) VDD
+ (18*K+tr) VDD  (19*K-tr) VDD
+ (19*K+tr) VDD  (20*K-tr) VDD
+ (20*K+tr) VDD  (21*K-tr) VDD
+ (21*K+tr) VDD  (22*K-tr) VDD
+ (22*K+tr) VDD  (23*K-tr) VDD
+ (23*K+tr) VDD  (24*K-tr) VDD
+ (24*K+tr) VDD  (25*K-tr) VDD
+ (25*K+tr) VDD  (26*K-tr) VDD
+ (26*K+tr) VDD  (27*K-tr) VDD
+ (27*K+tr) VDD  (28*K-tr) VDD
+ (28*K+tr) VDD  (29*K-tr) VDD
+ (29*K+tr) VDD  (30*K-tr) VDD
+ (30*K+tr) VDD  (31*K-tr) VDD
+ (31*K+tr) VDD  (32*K-tr) VDD
+ (32*K+tr) VDD  (33*K-tr) VDD
+ ]
vCLK (CLK 0) vsource type=pwl wave=\[
+ ( 0*K   ) 0    ( 1*K-tr) 0
+ ( 1*K+tr) 0    ( 2*K-tr) 0
+ ( 2*K+tr) VDD  ( 3*K-tr) VDD
+ ( 3*K+tr) VDD  ( 4*K-tr) VDD
+ ( 4*K+tr) 0    ( 5*K-tr) 0
+ ( 5*K+tr) 0    ( 6*K-tr) 0
+ ( 6*K+tr) VDD  ( 7*K-tr) VDD
+ ( 7*K+tr) VDD  ( 8*K-tr) VDD
+ ( 8*K+tr) 0    ( 9*K-tr) 0
+ ( 9*K+tr) 0    (10*K-tr) 0
+ (10*K+tr) VDD  (11*K-tr) VDD
+ (11*K+tr) VDD  (12*K-tr) VDD
+ (12*K+tr) 0    (13*K-tr) 0
+ (13*K+tr) 0    (14*K-tr) 0
+ (14*K+tr) VDD  (15*K-tr) VDD
+ (15*K+tr) VDD  (16*K-tr) VDD
+ (16*K+tr) 0    (17*K-tr) 0
+ (17*K+tr) 0    (18*K-tr) 0
+ (18*K+tr) VDD  (19*K-tr) VDD
+ (19*K+tr) VDD  (20*K-tr) VDD
+ (20*K+tr) 0    (21*K-tr) 0
+ (21*K+tr) 0    (22*K-tr) 0
+ (22*K+tr) VDD  (23*K-tr) VDD
+ (23*K+tr) VDD  (24*K-tr) VDD
+ (24*K+tr) 0    (25*K-tr) 0
+ (25*K+tr) 0    (26*K-tr) 0
+ (26*K+tr) VDD  (27*K-tr) VDD
+ (27*K+tr) VDD  (28*K-tr) VDD
+ (28*K+tr) 0    (29*K-tr) 0
+ (29*K+tr) 0    (30*K-tr) 0
+ (30*K+tr) VDD  (31*K-tr) VDD
+ (31*K+tr) VDD  (32*K-tr) VDD
+ (32*K+tr) 0    (33*K-tr) 0
+ ]
vD (D 0) vsource type=pwl wave=\[
+ ( 0*K   ) VDD  ( 1*K-tr) VDD
+ ( 1*K+tr) VDD  ( 2*K-tr) VDD
+ ( 2*K+tr) VDD  ( 3*K-tr) VDD
+ ( 3*K+tr) VDD  ( 4*K-tr) VDD
+ ( 4*K+tr) VDD  ( 5*K-tr) VDD
+ ( 5*K+tr) VDD  ( 6*K-tr) VDD
+ ( 6*K+tr) VDD  ( 7*K-tr) VDD
+ ( 7*K+tr) VDD  ( 8*K-tr) VDD
+ ( 8*K+tr) VDD  ( 9*K-tr) VDD
+ ( 9*K+tr) VDD  (10*K-tr) VDD
+ (10*K+tr) VDD  (11*K-tr) VDD
+ (11*K+tr) VDD  (12*K-tr) VDD
+ (12*K+tr) VDD  (13*K-tr) VDD
+ (13*K+tr) VDD  (14*K-tr) VDD
+ (14*K+tr) VDD  (15*K-tr) VDD
+ (15*K+tr) VDD  (16*K-tr) VDD
+ (16*K+tr) VDD  (17*K-tr) VDD
+ (17*K+tr) VDD  (18*K-tr) VDD
+ (18*K+tr) VDD  (19*K-tr) VDD
+ (19*K+tr) VDD  (20*K-tr) VDD
+ (20*K+tr) VDD  (21*K-tr) VDD
+ (21*K+tr) VDD  (22*K-tr) VDD
+ (22*K+tr) VDD  (23*K-tr) VDD
+ (23*K+tr) VDD  (24*K-tr) VDD
+ (24*K+tr) VDD  (25*K-tr) VDD
+ (25*K+tr) VDD  (26*K-tr) VDD
+ (26*K+tr) VDD  (27*K-tr) VDD
+ (27*K+tr) VDD  (28*K-tr) VDD
+ (28*K+tr) VDD  (29*K-tr) VDD
+ (29*K+tr) VDD  (30*K-tr) VDD
+ (30*K+tr) VDD  (31*K-tr) VDD
+ (31*K+tr) VDD  (32*K-tr) VDD
+ (32*K+tr) VDD  (33*K-tr) VDD
+ ]
\end{lstlisting}

\section{A.V \qquad Decoder SKILL code}

\paragraph{}
\begin{lstlisting}

load("./skill/Schematic.il")
procedure(DecoderSchematic(libName cellName WIDTH OUTPUTS)
  ;; Define the local variables
  prog((cvid x y pin pinName k h temp a b out inst)
    ;; Open the cell and its schematic view
    cvid = dbOpenCellViewByType(libName cellName "schematic" "schematic" "a")
	;; Clean the schematic view
    ece432DeleteObjectsSchematic(cvid)
	;; The coordinate of the origin point for the schematic
    x=0
    y=0	
    ;; Create top level pins	
	;; Create input pins
	for(u 0 WIDTH-1
		a=sprintf(nil "A%L" u)
		pin=ece432SchematicCreatePin(cvid a "input" x:y "R0")
      	y=y+0.5
    );; end for(u
	;; Create output pins
	for(q 0 OUTPUTS-1
		out=sprintf(nil "WL%L" q)
		pin=ece432SchematicCreatePin(cvid out "output" x:y "R0")
      	y=y+0.5
    );; end for(q
	;; Create inout pins
    foreach(pinName list("vdd" "gnd")
      pin=ece432SchematicCreatePin(cvid pinName "inputOutput" x:y "R0")
      y=y+0.5
    )
	;; The coordinate of the origin point for the schematic
    x=1
    y=0
	;; Create the decoder using tx-gates
	for(k 0 OUTPUTS-1
	  ;; Define the input and output net name of the first inverter
	  temp=k
	  if(mod(temp 2)==1 then
		add0="A0"
	  else
		add0="A0_"
      )
      temp=temp/2
	  if(mod(temp 2)==1 then
		add1="A1"
	  else
		add1="A1_"
      )
      temp=temp/2
	  if(mod(temp 2)==1 then
		add2="A2"
	  else
		add2="A2_"
      )
      temp=temp/2
	  if(mod(temp 2)==1 then
		add3="A3"
	  else
		add3="A3_"
      )
      temp=temp/2
	  if(mod(temp 2)==1 then
		add4="A4"
	  else
		add4="A4_"
      )
      temp=temp/2
	  if(mod(temp 2)==1 then
		add5="A5"
	  else
		add5="A5_"
      )
      temp=temp/2
	  if(mod(temp 2)==1 then
		add6="A6"
	  else
		add6="A6_"
      )
	  a=sprintf(nil "%s%s%s%s" add0 add1 add2 add3)
	  b=sprintf(nil "%s%s%s" add4 add5 add6)
	  out=sprintf(nil "WL%L" k)
	  ;; Instantiate the kth tx-gate with ece432SchematicCreateInst() procedure, 
	  ;; which is defined in the file 'ece432Schematic.il'. 
	  ;; This procedure needs 10 paramters
	  inst = ece432SchematicCreateInst(
		       ;; Cell view id
			   cvid 			   
			   ;; Name of the library containing the inverter cell
			   libName			   
		       ;; The tx-gate cell name  
	           "and_2"  			   
			   ;; Cell view type, always be "symbol" here
			   "symbol" 			   
			   ;; Instance name
		       sprintf(nil "AND_%L" k) 			   
			   ;; The following list defines the connections of the left side pins 
			   ;; in the symbol view; For our inverter, 'in' is the only pin on the left.
			   ;; The inverter pin 'in' is connected with the net 'inName'.
			   ;; If there are multiple pins on the left side, you can use the below way to define.
			   ;;    list(list("in1" in1net)
			   ;;         list("in2" in2net))
			   ;; If there is no pin on the left side, this parameter should be 'nil'.
	  		   list(list("a" a) 
			   		list("b" b))			   
			   ;; The following list defines the connections of the right side pins 
			   ;; in the symbol view; For our inverter, 'out' is the only pin on the right.
	  		   list(list("out" out)) 
			   ;; The following list defines the connections of the top side pins 
			   ;; in the symbol view; For our inverter, 'VDD' is the only pin on the top.
	  		   list(list("vdd"  "vdd")) 
			   ;; The following list defines the connections of the bottom side pins 
			   ;; in the symbol view; For our inverter, 'VSS' is the only pin on the bottom.
		       list(list("gnd"  "gnd"))		   
			   ;; Location of the instance
	           x:y 		   
			   ;; Rotation of the instance, such as "R0", "R90", "R180", "R270", "MX", "MY", ...
			   "R0"
			 );; end inst =       
	  y=y-2	  
	);; end for(k
	;; The coordinate of the origin point for the schematic
    x=-5
    y=0	
	;; Create the decoder using tx-gates
	for(k 0 15
	  ;; Define the input and output net name of the first inverter
	  temp=k
	  if(mod(temp 2)==1 then
		add0="A0"
	  else
		add0="A0_"
      )
      temp=temp/2
	  if(mod(temp 2)==1 then
		add1="A1"
	  else
		add1="A1_"
      )
      temp=temp/2
	  if(mod(temp 2)==1 then
		add2="A2"
	  else
		add2="A2_"
      )
      temp=temp/2
	  if(mod(temp 2)==1 then
		add3="A3"
	  else
		add3="A3_"
      )
	  a=sprintf(nil "%s" add0)
	  b=sprintf(nil "%s" add1)
	  c=sprintf(nil "%s" add2)
	  d=sprintf(nil "%s" add3)
	  out=sprintf(nil "%s%s%s%s" add0 add1 add2 add3)
	  ;; Instantiate the kth tx-gate with ece432SchematicCreateInst() procedure, 
	  ;; which is defined in the file 'ece432Schematic.il'. 
	  ;; This procedure needs 10 paramters
	  inst = ece432SchematicCreateInst(
		       ;; Cell view id
			   cvid 			   
			   ;; Name of the library containing the inverter cell
			   libName		   
		       ;; The tx-gate cell name  
	           "and_4"  		   
			   ;; Cell view type, always be "symbol" here
			   "symbol" 			   
			   ;; Instance name
		       sprintf(nil "AND4_%L" k) 			   
			   ;; The following list defines the connections of the left side pins 
			   ;; in the symbol view; For our inverter, 'in' is the only pin on the left.
			   ;; The inverter pin 'in' is connected with the net 'inName'.
			   ;; If there are multiple pins on the left side, you can use the below way to define.
			   ;;    list(list("in1" in1net)
			   ;;         list("in2" in2net))
			   ;; If there is no pin on the left side, this parameter should be 'nil'.
	  		   list(list("a" a) 
			   		list("b" b)
			   		list("c" c)
			   		list("d" d))		   
			   ;; The following list defines the connections of the right side pins 
			   ;; in the symbol view; For our inverter, 'out' is the only pin on the right.
	  		   list(list("out" out)) 
			   ;; The following list defines the connections of the top side pins 
			   ;; in the symbol view; For our inverter, 'VDD' is the only pin on the top.
	  		   list(list("vdd"  "vdd")) 
			   ;; The following list defines the connections of the bottom side pins 
			   ;; in the symbol view; For our inverter, 'VSS' is the only pin on the bottom.
		       list(list("gnd"  "gnd"))		   
			   ;; Location of the instance
	           x:y 		   
			   ;; Rotation of the instance, such as "R0", "R90", "R180", "R270", "MX", "MY", ...
			   "R0"
			 );; end inst =	       
	  y=y-2  
	);; end for(k
	;; The coordinate of the origin point for the schematic
    x=-8
    y=0	
	;; Create the decoder using tx-gates
	for(k 0 7
	  ;; Define the input and output net name of the first inverter
	  temp=k
	  if(mod(temp 2)==1 then
		add0="A4"
	  else
		add0="A4_"
      )
      temp=temp/2
	  if(mod(temp 2)==1 then
		add1="A5"
	  else
		add1="A5_"
      )
      temp=temp/2
	  if(mod(temp 2)==1 then
		add2="A6"
	  else
		add2="A6_"
      )
	  a=sprintf(nil "%s" add0)
	  b=sprintf(nil "%s" add1)
	  c=sprintf(nil "%s" add2)
	  out=sprintf(nil "%s%s%s" add0 add1 add2)
	  ;; Instantiate the kth tx-gate with ece432SchematicCreateInst() procedure, 
	  ;; which is defined in the file 'ece432Schematic.il'. 
	  ;; This procedure needs 10 paramters
	  inst = ece432SchematicCreateInst(
		       ;; Cell view id
			   cvid 		   
			   ;; Name of the library containing the inverter cell
			   libName		   
		       ;; The tx-gate cell name  
	           "and_3"  	   
			   ;; Cell view type, always be "symbol" here
			   "symbol" 	   
			   ;; Instance name
		       sprintf(nil "AND3_%L" k) 	   
			   ;; The following list defines the connections of the left side pins 
			   ;; in the symbol view; For our inverter, 'in' is the only pin on the left.
			   ;; The inverter pin 'in' is connected with the net 'inName'.
			   ;; If there are multiple pins on the left side, you can use the below way to define.
			   ;;    list(list("in1" in1net)
			   ;;         list("in2" in2net))
			   ;; If there is no pin on the left side, this parameter should be 'nil'.
	  		   list(list("a" a) 
			   		list("b" b)
			   		list("c" c))	   
			   ;; The following list defines the connections of the right side pins 
			   ;; in the symbol view; For our inverter, 'out' is the only pin on the right.
	  		   list(list("out" out)) 
			   ;; The following list defines the connections of the top side pins 
			   ;; in the symbol view; For our inverter, 'VDD' is the only pin on the top.
	  		   list(list("vdd"  "vdd")) 
			   ;; The following list defines the connections of the bottom side pins 
			   ;; in the symbol view; For our inverter, 'VSS' is the only pin on the bottom.
		       list(list("gnd"  "gnd"))   
			   ;; Location of the instance
	           x:y    
			   ;; Rotation of the instance, such as "R0", "R90", "R180", "R270", "MX", "MY", ...
			   "R0"
			 );; end inst =
	  y=y-2
	);; end for(k
;; ******************************Address inverters***************************
	;; The coordinate of the origin point for the schematic
    x=-12
    y=0	
	;; Create the decoder using tx-gates
	for(k 0 6
	  in=sprintf(nil "A%L" k)
	  out=sprintf(nil "A%L_" k)
	  ;; Instantiate the kth tx-gate with ece432SchematicCreateInst() procedure, 
	  ;; which is defined in the file 'ece432Schematic.il'. 
	  ;; This procedure needs 10 paramters
	  inst = ece432SchematicCreateInst(
		       ;; Cell view id
			   cvid 			   
			   ;; Name of the library containing the inverter cell
			   libName			   
		       ;; The tx-gate cell name  
	           "inv"  			   
			   ;; Cell view type, always be "symbol" here
			   "symbol" 			   
			   ;; Instance name
		       sprintf(nil "INVERTER_%L" k) 			   
			   ;; The following list defines the connections of the left side pins 
			   ;; in the symbol view; For our inverter, 'in' is the only pin on the left.
			   ;; The inverter pin 'in' is connected with the net 'inName'.
			   ;; If there are multiple pins on the left side, you can use the below way to define.
			   ;;    list(list("in1" in1net)
			   ;;         list("in2" in2net))
			   ;; If there is no pin on the left side, this parameter should be 'nil'.
	  		   list(list("in" in))			   
			   ;; The following list defines the connections of the right side pins 
			   ;; in the symbol view; For our inverter, 'out' is the only pin on the right.
	  		   list(list("out" out)) 
			   ;; The following list defines the connections of the top side pins 
			   ;; in the symbol view; For our inverter, 'VDD' is the only pin on the top.
	  		   list(list("vdd"  "vdd")) 
			   ;; The following list defines the connections of the bottom side pins 
			   ;; in the symbol view; For our inverter, 'VSS' is the only pin on the bottom.
		       list(list("gnd"  "gnd"))		   
			   ;; Location of the instance
	           x:y 			   
			   ;; Rotation of the instance, such as "R0", "R90", "R180", "R270", "MX", "MY", ...
			   "R0"
			 );; end inst =	       
	  y=y-2	  
	);; end for(k
	;; Check, save and close the cell view
	schCheck(cvid)  ;; check the schematic connectivity
    dbSave(cvid)
    dbClose(cvid)	
    return(t)	
  );; end prog 
);; end procedure
/*====================================================================
 'let' is the main entrance of the skill program.
====================================================================*/
let((WIDTH libName cellName)  ;; N, libName and cellName are variables
  WIDTH = 7
  OUTPUTS = 74
  ;; The name of the library we will put the new cell in
  libName="sumitfpga"
  sprintf(cellName "RowDecoder")
  DecoderSchematic(libName cellName WIDTH OUTPUTS)
  printf("=== Cell %L Schematic has been created! ===\n", cellName)
) ;;end let
\end{lstlisting}
 % filename in curly brackets

\clearemptydoublepage
%
% Add index
\printindex
%
\nocite{*}
\end{document} 
